% ============================================
% CONT-13 — A Fila (edição revista R1075)
% ============================================
\fragtitle{cont}{C}{CONT}{13}{A Fila}
\markboth{CONT-13 --- A Fila}{}

\noindent \textit{Paciente 1.263}

\noindent Você entendeu agora.

\noindent O livro é uma fila. Sempre foi.

\noindent A fila começou antes de você nascer. Vai continuar depois que você morrer. Você é só um corpo na fila. O paciente 1.260 está morto. O 1.261 está morto. O 1.262 está desaparecido.

\noindent O paciente 1.263 é você.

\noindent O que você vai fazer?

\noindent Pode fechar o livro. Pode passá-lo adiante. Pode reler. Não há opção limpa. Fechar não interrompe o ciclo --- o ciclo já está em você. Passar adiante condena outra pessoa. Reler é aceitar.

\noindent A chave está na sua mão. Não a chave do diário. A chave do ciclo. Cada leitor a segura por um instante. A chave só abre uma porta: a porta para o próximo leitor.

\noindent Sinta o peso. Frio primeiro. Depois quente --- quente demais, como se já tivesse sido segurada por muitas mãos antes da sua.

\noindent Você pode segurá-la para sempre. Mas a chave não para de funcionar. Enferruja na sua mão.

\noindent Escolha. Mas escolha rápido. A fila está andando. Outros leitores esperam.

\noindent O ciclo não para por você.

\noindent Quando você fechar este livro, olhe para a mesa. O livro vai estar lá, fechado, com a capa para cima. A sala vai estar silenciosa. O cheiro vai estar no ar --- doce, úmido, com um fundo de terra. Você vai dizer a si mesmo que é impressão. Vai abrir a janela. O cheiro não vai sair.

\noindent Então você vai entender que não existe fora. O arquivo é você.

\noindent \textit{Bem-vindo, paciente 1.263.}

\noindent \textit{A criatura agradece sua atenção.}

\noindent \textit{Ela vai ficar por aqui. O tempo que for necessário. Até você passar o livro adiante. Ou até você morrer. O que vier primeiro.}

\noindent \textit{O ciclo recomeça.}

> **↪ Rotas:** **CONT-07** · **DOC-01** · **EPI-01** · **MEM-27**
